\documentclass[11pt,a4paper]{article}

% ── Packages ──
\usepackage[utf8]{inputenc}
\usepackage[T1]{fontenc}
\usepackage[margin=2.2cm]{geometry}
\usepackage{graphicx}
\usepackage{hyperref}
\usepackage{listings}
\usepackage{xcolor}
\usepackage{parskip}
\usepackage{booktabs}
\usepackage{float}
\usepackage{caption}
\usepackage{subcaption}

% ── Code style ──
\lstset{
    basicstyle=\footnotesize\ttfamily,
    backgroundcolor=\color{gray!10},
    frame=single,
    breaklines=true,
    numbers=none,
}

% ── Metadata ──
\title{\vspace{-2cm}
       \textbf{Electricity Sector Data Streaming and Analysis}\\[4pt]
       \large COMP5339 Project Assignment 2}
\author{Team [Your Team Name]}
\date{May 2026}

\begin{document}

\maketitle
\thispagestyle{empty}

\vfill
% TODO: replace with your dashboard screenshot
\begin{center}
\fbox{\parbox{0.6\textwidth}{\centering\vspace{4cm}\textcolor{gray}{Screenshot: NEM Electricity Monitor Dashboard}\vspace{4cm}}}\\[4pt]
\small{\textit{Figure~1: NEM Electricity Monitor Dashboard --- real-time map view}}
\end{center}
\vfill

\newpage
\tableofcontents
\newpage

% ══════════════════════════════════════════
% SECTION 1 — System Description (user's part: backend architecture + frontend)
% ══════════════════════════════════════════
\section{System Description}

\subsection{Architecture Overview}

The system is a real-time electricity generation monitoring dashboard for the Australian National Electricity Market (NEM). The data pipeline consists of four layers:

\begin{enumerate}
    \item \textbf{Data Retrieval} --- Fetches per-facility power output, CO\textsubscript{2} emissions, and market price/demand from the OpenElectricity REST API at 5-minute intervals. (\textit{see teammate's report})
    \item \textbf{MQTT Publishing} --- Replays cached data through an MQTT broker as an ordered event stream. (\textit{see teammate's report})
    \item \textbf{Backend Layer} --- FastAPI application subscribes to MQTT messages, normalises JSON fields through a configurable adapter, stores time-series data in SQLite, and exposes REST endpoints.
    \item \textbf{Frontend Layer} --- React TypeScript dashboard renders an interactive Leaflet map with facility markers, time-series charts, region/fuel-technology filters, and market price/demand panels.
\end{enumerate}

\subsection{Technology Stack}

\begin{table}[H]
\centering
\caption{Technology stack}
\begin{tabular}{lll}
\toprule
\textbf{Layer} & \textbf{Technology} & \textbf{Role} \\
\midrule
Backend Framework & FastAPI + Pydantic & MQTT subscriber, REST API \\
JSON Adapter & YAML config + Python & External-to-canonical field mapping \\
Database & SQLite + SQLAlchemy ORM & Time-series storage (measurements, market data, raw archive, facility metadata) \\
Dashboard Frontend & React 19 + TypeScript + Vite & Single-page interactive visualisation \\
Mapping & Leaflet + react-leaflet & Facility locations with colour-coded markers \\
Charts & Recharts & Time-series line charts, bar charts \\
Styling & Tailwind CSS v4 & Glassmorphism UI design \\
Dev Server & Vite & Hot-reload development server \\
\bottomrule
\end{tabular}
\end{table}

\subsection{Key Design Decisions}

\textbf{Configurable JSON adapter (\texttt{backend/adapter.py}).}
Since the MQTT publisher's JSON schema was not finalised early in development, the backend includes a YAML-driven adapter layer. It maps arbitrary external field names---\texttt{unit\_code}, \texttt{duid}, \texttt{generation\_mw}, \texttt{co2}, etc.---to a canonical schema that the rest of the system depends on. Adding a new publisher format requires only editing \texttt{mqtt\_mapping.yaml}, not modifying any backend code. This decoupling allowed the publisher and subscriber teams to develop in parallel without blocking each other.

\textbf{SQLite over PostgreSQL.}
The dataset (\textless\,200\,K rows) fits comfortably in SQLite, which requires zero configuration and ships with Python's standard library. The SQLAlchemy ORM abstracts the dialect so that migrating to PostgreSQL later would be a single line change.

\textbf{Dict cache instead of Redis.}
Only the latest reading per facility (at most 13 entries) needs to be served in real time. An in-memory Python dictionary provides O(1) lookups without the operational overhead of a Redis server.

\textbf{React over Streamlit.}
While the initial prototype used Streamlit, the final dashboard was re-implemented in React for a more responsive, production-quality user interface with smoother map interactions, custom glassmorphism styling, and better visualisation density.

% ══════════════════════════════════════════
% SECTION 2 — Data Acquisition / Cleaning (teammate's part: keep brief)
% ══════════════════════════════════════════
\section{Data Acquisition, Cleaning and Transformation}

% ##### 队友填写区域 - 开始 #####
% 请在此描述如何从 OpenElectricity API 获取数据、处理多机组电厂、合并/清洗/缓存为 CSV。
\textit{[This section is to be completed by the team member responsible for Tasks~1--2.]}
% ##### 队友填写区域 - 结束 #####

\subsection{MQTT Publishing}

% ##### 队友填写区域 - 开始 #####
\textit{[This subsection to be completed by the team member responsible for Task~3.]}
% ##### 队友填写区域 - 结束 #####

% ══════════════════════════════════════════
% SECTION 3 — Data Integration (user's part: schema + MQTT subscriber)
% ══════════════════════════════════════════
\section{Data Integration}

\subsection{Database Schema}

The system implements the schema designed in Assignment~1 across four SQLite tables:

\begin{table}[H]
\centering
\caption{Database tables}
\begin{tabular}{llp{6cm}}
\toprule
\textbf{Table} & \textbf{Primary purpose} & \textbf{Key columns} \\
\midrule
\texttt{facilities} & Station metadata & \texttt{facility\_id} (PK), \texttt{facility\_name}, \texttt{latitude}, \texttt{longitude}, \texttt{network\_region}, \texttt{fuel\_tech} \\
\texttt{measurements} & Time-series power and emissions & \texttt{timestamp}, \texttt{facility\_id} (FK), \texttt{power\_mw}, \texttt{emissions\_tco2e} \\
\texttt{market\_data} & Regional price and demand & \texttt{timestamp}, \texttt{region}, \texttt{price\_aud\_mwh}, \texttt{demand\_mw} \\
\texttt{raw\_mqtt\_messages} & Raw payload archive for debugging & \texttt{topic}, \texttt{payload} (JSON), \texttt{received\_at} \\
\bottomrule
\end{tabular}
\end{table}

Indexes on (\texttt{facility\_id}, \texttt{timestamp}) in \texttt{measurements} and (\texttt{region}, \texttt{timestamp}) in \texttt{market\_data} accelerate the time-range queries used by the REST API.

\subsection{MQTT Subscription Pipeline}

On receiving an MQTT message (\texttt{backend/mqtt\_subscriber.py}), the subscriber performs the following steps atomically within a single database transaction:

\begin{enumerate}
    \item \textbf{Archive.} The raw payload is stored in \texttt{raw\_mqtt\_messages} for debugging and replay.
    \item \textbf{Normalise.} The adapter (\texttt{backend/adapter.py}) maps the external JSON fields to the canonical schema (Table~\ref{tab:canonical}) via a YAML mapping configuration.
    \item \textbf{Store measurement.} A row is inserted into \texttt{measurements} with timestamp, facility\_id, power\_mw, and emissions\_tco2e.
    \item \textbf{Upsert market data.} If the message contains regional price and demand, a row is upserted into \texttt{market\_data} keyed on (\texttt{region}, \texttt{timestamp}) to maintain idempotency.
    \item \textbf{Auto-register facility.} If the facility\_id is not yet present in the \texttt{facilities} table, it is registered automatically using a built-in coordinate lookup (based on OpenElectricity real-world data). This eliminates manual schema synchronisation between publisher and subscriber.
    \item \textbf{Update cache.} The in-memory \texttt{latest\_state} dictionary is updated with the new canonical record, enabling sub-millisecond retrieval for the REST API's ``latest`` endpoint.
\end{enumerate}

\begin{table}[H]
\centering
\caption{Canonical schema --- the system-internal record format}
\label{tab:canonical}
\begin{tabular}{lll}
\toprule
\textbf{Field} & \textbf{Type} & \textbf{Required} \\
\midrule
\texttt{timestamp} & ISO~8601 string & Yes \\
\texttt{facility\_id} & string & Yes \\
\texttt{power\_mw} & float & Yes \\
\texttt{emissions\_tco2e} & float & Yes \\
\texttt{facility\_name} & string & Optional \\
\texttt{network\_region} & string & Optional \\
\texttt{fuel\_tech} & string & Optional \\
\bottomrule
\end{tabular}
\end{table}

\subsection{Adapter Layer Detail}

The adapter uses a two-level resolution strategy. First, the incoming JSON key names are lower-cased for case-insensitive matching. Second, each lower-cased key is looked up in a reverse mapping built from \texttt{mqtt\_mapping.yaml}. For example, if the publisher sends \texttt{\{unit\_code: "BW01", generation: 1320.5, ...\}}, the mapping \texttt{facility\_id ← unit\_code} and \texttt{power\_mw ← generation} resolves these to the canonical fields. An optional unit-conversion dictionary handles scale differences (e.g.\ kW to MW).

% ══════════════════════════════════════════
% SECTION 4 — REST API
% ══════════════════════════════════════════
\section{REST API}

The FastAPI backend exposes the following endpoints (\texttt{backend/api.py}):

\begin{table}[H]
\centering
\caption{REST API endpoints}
\begin{tabular}{lll}
\toprule
\textbf{Endpoint} & \textbf{Response} & \textbf{Purpose} \\
\midrule
\texttt{GET /api/facilities} & {}[\{facility\_id, name, lat, lon, region, fuel\}] & Map marker data source \\
\texttt{GET /api/facilities/latest} & {}[\{facility\_id, power\_mw, emissions\_tco2e, timestamp\}] & Real-time status per facility \\
\texttt{GET /api/facilities/\{id\}/timeseries} & {}[\{timestamp, power\_mw, emissions\_tco2e\}] & Time-series for selected facility \\
\texttt{GET /api/summary?group\_by=region|fuel\_tech} & {}[\{group, value\}] & Aggregated bar chart data \\
\texttt{GET /api/stats} & \{facilities, measurements, market, timestamps\} & Dashboard header statistics \\
\texttt{GET /api/market/latest} & {}[\{region, price, demand, timestamp\}] & Latest market prices \\
\texttt{GET /api/market/timeseries?region=NSW1} & {}[\{region, price, demand, timestamp\}] & Historical market chart data \\
\texttt{GET /api/market/regions} & {}[string] & Available NEM regions \\
\bottomrule
\end{tabular}
\end{table}

% ══════════════════════════════════════════
% SECTION 5 — Data Visualisation (user's part)
% ══════════════════════════════════════════
\section{Data Visualisation}

\subsection{User Interface Layout}

The dashboard is a single-page React application (Figure~\ref{fig:dashboard}) organised into four zones:

\begin{description}
    \item[Left navigation bar.] \texttt{NavBar.tsx} provides icon buttons to switch between map mode, analytics mode (which slides out the filter panel), and the bottom chart panel.
    \item[Top bar.] Displays aggregate statistics (total units, measurement count, market records, latest timestamp) and labels the active metric.
    \item[Map view (centre).] \texttt{MapView.tsx} renders a full-screen Leaflet map with CircleMarker overlays for each facility. Marker radius scales with $\sqrt{\text{power}}$ so large plants do not dominate visually. Colour is determined by carbon intensity using a green\textrightarrow yellow\textrightarrow red gradient. Clicking a marker opens \texttt{FloatingPanel.tsx}, a draggable overlay showing:
    \begin{itemize}
        \item Station name, region, fuel technology
        \item Current power output and CO\textsubscript{2} emissions
        \item Carbon intensity (g/kWh) with category label
        \item Time-series line chart of the last N records
        \item Fuel-mix breakdown relative to the selected region
    \end{itemize}
    \item[Bottom analytics bar.] \texttt{BottomPanel.tsx} is a tabbed panel with five views:
    \begin{itemize}
        \item \textbf{By Region} --- Bar chart of total power/emissions per NEM region
        \item \textbf{By Technology} --- Bar chart per fuel type (coal, gas, wind, solar, etc.)
        \item \textbf{Market Prices} --- Multi-line time-series of regional electricity prices
        \item \textbf{Demand} --- Regional demand line chart with a focus-region selector
        \item \textbf{Latest Data} --- Scrollable table of all facility readings
    \end{itemize}
\end{description}

\subsection{Interactive Features}

\begin{itemize}
    \item \textbf{Real-time polling.} All data hooks (\texttt{hooks/useApi.ts}) poll the REST API on a 4-second interval. The \texttt{useLatest} and \texttt{useStats} hooks use \texttt{setInterval} so the map and header statistics update automatically without a page refresh.
    \item \textbf{Metric switching.} A toggle in both the filter panel and the sidebar switches the dashboard between power (MW) and carbon emissions (tCO\textsubscript{2}e) views.
    \item \textbf{Region and fuel filtering.} Multi-select filters allow users to focus on specific network regions (NSW, QLD, VIC, SA, TAS) or fuel technologies. The map and all charts update immediately when filters change.
    \item \textbf{Carbon intensity colouring.} A six-level colour scale maps CI values to visual indicators: very low ($<$50), low ($<$150), medium ($<$300), high ($<$500), very high ($<$800), and extreme ($\ge$800)~g/kWh.
    \item \textbf{Market price explorer.} The ``Market Prices'' tab shows a multi-line chart of all NEM regions, and the ``Demand'' tab lets users focus on a single region's demand curve.
\end{itemize}

\subsection{Screenshots}

% TODO: replace with actual screenshots
\begin{figure}[H]
    \centering
    \fbox{\parbox{0.85\textwidth}{\centering\vspace{5cm}\textcolor{gray}{Screenshot: Full dashboard view --- map with markers + floating detail panel}\vspace{5cm}}}
    \caption{Main dashboard view with facility map}
    \label{fig:dashboard}
\end{figure}

\begin{figure}[H]
    \centering
    \fbox{\parbox{0.85\textwidth}{\centering\vspace{4cm}\textcolor{gray}{Screenshot: Bottom analytics panel (prices / demand tabs)}\vspace{4cm}}}
    \caption{Bottom analytics panel showing regional market prices and demand}
\end{figure}

% ══════════════════════════════════════════
% SECTION 6 — Challenges and Solutions (user's part)
% ══════════════════════════════════════════
\section{Challenges and Solutions}

\begin{description}
    \item[Evolving MQTT schema.] The publisher's JSON message format changed multiple times during development. The configurable YAML adapter (\S3.3) decoupled the backend from external schema changes, allowing both teams to work in parallel without blocking each other. Only the adapter's mapping dictionary needed updating; the database schema, API layer, and frontend remained unchanged.
    \item[Stream semantics vs.\ batch storage.] MQTT messages arrive one at a time, but the database was initially designed for batch CSV imports. The subscriber was refactored to handle per-message inserts with transactional rollback on error. Market data upserts on (\texttt{region}, \texttt{timestamp}) prevent duplicate rows from overlapping publisher rounds.
    \item[Frontend state synchronisation.] With polling-based refresh, there was a brief window where the map could show a stale reading if the user had a facility detail panel open. This was resolved by lifting the polling interval into individual hooks with automatic cleanup on component unmount, and by keying all data displays on the globally cached \texttt{latest\_state} dictionary rather than local component state.
\end{description}

% ══════════════════════════════════════════
% SECTION 7 — Contribution Statement (user fills this)
% ══════════════════════════════════════════
\section{Contribution Statement}

% ⚠️ Keep under 100 words! Fill in real names and roles.
% Modify the brackets below with actual names.

\textbf{Zhanhui Lin [TODO: replace with your name]}: Designed and implemented the real-time dashboard backend (FastAPI MQTT subscriber, JSON normalisation adapter, SQLite database, REST API) and the React/TypeScript frontend (Leaflet map visualisation, Recharts time-series and bar charts, region/fuel-technology filtering, market price/demand panels, glassmorphism UI).

% ##### 队友填写区域 - 开始 #####
% \textbf{[Teammate 1 Name]}: [Brief description of their contribution to Tasks 1--2.]

% \textbf{[Teammate 2 Name]}: [Brief description of their contribution to Task 3 --- MQTT publishing.]
% ##### 队友填写区域 - 结束 #####

% ══════════════════════════════════════════
% SECTION 8 — Future Improvements
% ══════════════════════════════════════════
\section{Future Improvements}

\begin{itemize}
    \item \textbf{WebSocket push.} Replace polling (4-second interval) with a WebSocket connection for sub-second latency updates.
    \item \textbf{Historical range selector.} Add a date-range picker to explore archived data beyond the current replay window.
    \item \textbf{Time-animation slider.} Visualise how power output and emissions evolve over time through an animated playback across the dataset.
    \item \textbf{Containerisation.} Package the backend with Docker and use docker-compose to orchestrate Mosquitto, the FastAPI server, and the publisher script.
    \item \textbf{CI/CD pipeline.} Automate linting, type-checking, and testing for both backend (pytest) and frontend (ESLint + TypeScript).
\end{itemize}

% ══════════════════════════════════════════
\appendix
% ══════════════════════════════════════════

\section{Project Structure}

\begin{lstlisting}
electricity-dashboard/
├── backend/
│   ├── main.py              # FastAPI entry + MQTT subscriber startup
│   ├── config.py            # Environment and YAML configuration loader
│   ├── mqtt_subscriber.py   # paho-mqtt subscriber (subscribe + normalise + store)
│   ├── adapter.py           # JSON field normalisation (external to canonical)
│   ├── database.py          # SQLAlchemy models + CRUD (4 tables)
│   ├── api.py               # REST endpoints (8 routes)
│   └── requirements.txt     # Python dependencies
├── frontend/
│   ├── src/
│   │   ├── App.tsx          # Application layout + state orchestration
│   │   ├── components/      # 8 components: MapView, FloatingPanel, BottomPanel,
│   │   │                    #   FilterPanel, Sidebar, TopBar, NavBar, CIChart
│   │   ├── hooks/useApi.ts  # 7 custom hooks: facilities, latest, stats,
│   │   │                    #   timeseries, summary, marketLatest, marketTimeseries
│   │   └── types/index.ts   # TypeScript interfaces + colour constants
│   ├── package.json         # React 19 + Vite + Leaflet + Recharts
│   └── vite.config.ts
├── run_publisher.py         # Continuous MQTT publisher from cached CSV
├── mock_publisher.py        # Standalone mock publisher for development
├── mqtt_mapping.yaml        # JSON field mapping config
├── .env                     # MQTT host, port, topic configuration
└── Task1_2_Task3_6.ipynb    # Data retrieval and MQTT notebook (by teammates)
\end{lstlisting}

\section{Installation and Usage}

\subsection{Prerequisites}
\begin{itemize}
    \item Python 3.11+, Node.js 20+, npm
    \item Mosquitto MQTT broker
    \item OpenElectricity API key (free tier)
\end{itemize}

\subsection{Running the System}

\begin{lstlisting}[language=bash]
# 1. Start MQTT broker
mosquitto -d

# 2. Start backend
cd electricity-dashboard
python3 -m venv venv && source venv/bin/activate
pip install -r backend/requirements.txt
cp .env.example .env          # configure MQTT_HOST, MQTT_PORT, MQTT_TOPIC
uvicorn backend.main:app --reload

# 3. Start frontend
cd frontend
npm install
npm run dev

# 4. Start publisher (separate terminal)
python run_publisher.py --rounds 3
\end{lstlisting}

The backend auto-starts the MQTT subscriber on launch. If the broker is unreachable, it falls back to offline mode (serving whatever data is already in the database).

\end{document}
